\chapter{Введение: постановка задачи пятого тома}

Четвёртый том сохранил воспроизводимое математическое ядро, но завершился
оценкой научной готовности \(R_{\rm sci}=4/10\) и
\begin{equation}N_{\text{физически замкнутых}}=0.\end{equation}
Ни один наблюдаемый сектор не был полностью выведен из единой исходной
конструкции. Отдельные представления, действия, нормировки и слепые проверки
работали, но не соединялись одной мерой и одним вариационным принципом.

\section{Почему необходима новая версия}

Повторяющееся препятствие состояло не в нехватке численного совпадения, а в
отсутствии заранее заданной архитектуры, которая одновременно фиксирует
пространство состояний, меру, вакуум, относительные веса секторов,
калибровочный фактор, допустимые контрчлены и отображение к наблюдаемым.

Можно было построить дефектное ядро, поколенческий проектор или отдельный
спектральный коэффициент, но конденсат, сектор фиксированного заряда, знак
квартичного функционала либо нормировку приходилось выбирать отдельно.

\begin{definition}[Родительская архитектура]
Родительской архитектурой называется совокупность
\begin{equation}
 \mathfrak P=
 (\mathcal A,\mathcal H,D,J,\gamma;
  \mathfrak C,\mathfrak S,\mu;
  \mathcal F,\mathcal Q,\mathcal R),
\end{equation}
где первая группа задаёт операторные и конечно-геометрические данные,
\(\mathfrak C\) --- класс носителей или граничных фонов,
\(\mathfrak S\) --- пространство состояний, \(\mu\) --- меру и статистику,
\(\mathcal F\) --- один исходный функционал, \(\mathcal Q\) --- физический
фактор по калибровочным и BV-симметриям, а \(\mathcal R\) --- заранее
фиксированное отображение к обучающим и слепым наблюдаемым.
\end{definition}

Эта структура шире спектральной тройки: предыдущие тома показали, что один
спектр не определяет ансамбль, веса топологий, конечные контрчлены и
физический фактор.

\section{Задача общего вариационного источника}

\begin{definition}[Общий вариационный источник]
Архитектура считается допустимой к физической проверке, если носитель,
граничный сектор или конечная геометрия заданы до обращения к наблюдаемым;
исходный функционал фиксирован заранее; мера охватывает все вспомогательные
секторы; физическая стационарная точка существует после полного
факторизования; нулевые моды совпадают с доказанными калибровочными
направлениями; условия перенормировки заданы до слепых проверок; одна и та же
нормировка управляет не менее чем двумя независимыми секторами.
\end{definition}

\section{Что наследуется из предыдущих томов}

Сохраняются только доказанные модули: спектры \(\mathbb{RP}^3\) и накрытия,
минимальная конечная алгебра Стандартной модели, нормированные тепловые
состояния, относительные селекторы Пати--Салама, тетраэдрическая
поколенческая геометрия, голономия \(2\pi/3\), одна мода Майораны,
вильсоновские коэффициенты и операторные тождества гессиана. Их наличие не
считается доказательством общего действия.

\section{Запрещённый багаж}

\begin{nogocriterion}[Запрет повторного использования закрытых механизмов]
Без нового родительского вывода запрещается подбирать пороговые массы,
выбирать отображения Юкавы по качеству согласия с массами или CKM,
объявлять обычный \(\Tr D_F^4\) источником поколенческого отображения
момента с противоположным знаком, восстанавливать удалённые лишние 2-формы
на особом фоне, использовать мнимое преобразование Хаббарда--Стратоновича
без исходного фермионного ядра, сравнивать топологии со свободными
кривизными весами и выводить абсолютный масштаб из неконтролируемого
локального разложения теплового ядра.
\end{nogocriterion}

\section{Три допустимых класса архитектуры}

\subsection{Совместная вариация состояния и носителя}

Первый класс использует нормированное состояние \(\rho\) и геометрию
носителя \(M\) как совместные переменные:
\begin{equation}
 \mathcal F_{\rm C}[\rho,M]
 =\Tr(\rho H_C[M])+\Gamma_{\text{фл}}[\rho,M]-\tau^{-1}S(\rho),
 \qquad H_C=-\tau^{-1}\log\widehat C.
\end{equation}
Главные вопросы --- знак исходного функционала, мера на носителях и конечные
гравитационно-топологические веса.

\subsection{Архитектура граничного источника}

Второй класс строит единый граничный интеграл по траекториям или
суперсвязность, действующие на стерильной корневой линии, поле спаривания
заряда два, вещественный поколенческий триплет и квантованный ротор. Один
след должен породить фиксированный заряд, конденсат, тетраэдрическую ось,
голономию и ровно одну моду Майораны.

\subsection{Новая конечная геометрия}

Третий класс содержательно меняет конечную геометрию: допускает носитель
ранга не меньше трёх, несколько некоммутирующих проекторов, производный
четырёхфермионный канал, изменённое дифференциальное исчисление или
динамическое кручение Картана. Простое добавление ещё одного свободного поля
не считается новой архитектурой.

\section{Проверка V.0 до феноменологии}

\begin{protocol}[Проверка V.0: заморозка архитектуры]
Каждый класс обязан представить один исходный объект, один функционал без
свободных межсекторных весов, полный реестр алгебры, пространства, меры и
статистики, список допустимых контрчленов, одну вычислимую вне массовой
оболочки проверку гессиана, ранга или знака и заранее записанный двоичный
критерий. После выбора одного класса остальные не могут передавать ему свои
поля и коэффициенты.
\end{protocol}

\section{Критерий завершённости Тома V}

\begin{definition}[Математическое замыкание]
Все переменные, функционал, правила контрчленов, физический фактор и полный
гессиан определены и воспроизводимы до подстановки наблюдаемых данных.
\end{definition}

\begin{definition}[Межсекторное замыкание]
Одна мера и одно соглашение о следе фиксируют не менее двух независимых
чувствительных к нормировке секторных весов.
\end{definition}

\begin{definition}[Физическое замыкание]
После одного размерного обучающего входа полный ход группы перенормировки и
порогов проходит две заранее зарегистрированные безразмерные слепые проверки
из независимых секторов без дополнительной подгонки.
\end{definition}

\begin{protocol}[Критерий завершённости]
Требуются: выбор одной архитектуры до феноменологии; положительная проверка
полного гессиана и фактора BV/BRST; единая межсекторная нормировка;
пороговый расчёт без масс, выведенных из цели; независимое воспроизведение;
сохранение всех неудачных наблюдаемых и отрицательных мод в итоговом отчёте.
\end{protocol}

\section{Условие ранней остановки}

\begin{nogocriterion}[Ранняя остановка]
Если ни один класс не фиксирует исходный функционал и физический гессиан до
наблюдаемых входов, Том V завершается классификацией препятствий и не
открывает новые подгонки масс, смешиваний, связей или абсолютного масштаба.
\end{nogocriterion}

Первый рабочий шаг --- сравнительная проверка трёх архитектур без
наблюдаемых чисел.